\documentclass[11pt]{article}
% Shared preamble for Theseus user guides.
% Compiled with pdflatex. See docs/guides/BUILD.md for rationale.

\usepackage[a4paper,margin=0.85in]{geometry}
\usepackage[T1]{fontenc}
\usepackage[utf8]{inputenc}

% Body: Palatino (via mathpazo). Headings: Helvetica (sans).
% Palatino is requested in the house style as a substitute for Charter
% because Charter is not part of the standard TeX Live distribution
% that the build container ships.
\usepackage{mathpazo}
\usepackage[scaled=0.92]{helvet}
\usepackage{courier}
\renewcommand{\familydefault}{\rmdefault}

\usepackage{microtype}
\usepackage{parskip}
\usepackage{enumitem}
\usepackage{booktabs}
\usepackage{longtable}
\usepackage{titlesec}
\usepackage{xcolor}
\usepackage{fancyhdr}
\usepackage{fancyvrb}
\usepackage{graphicx}
% Allow `_` in text mode without escaping. The screenshot file names and
% some environment-variable references in the guides contain underscores.
\usepackage{underscore}
\usepackage{tcolorbox}
\tcbuselibrary{breakable,skins}
\usepackage[hidelinks,colorlinks=true,linkcolor=black,urlcolor=teal,citecolor=black]{hyperref}

% Sans for chapter and section headings; serif body keeps prose readable
% on screen and in print.
\titleformat{\section}{\sffamily\Large\bfseries}{\thesection}{0.7em}{}
\titleformat{\subsection}{\sffamily\large\bfseries}{\thesubsection}{0.6em}{}
\titleformat{\subsubsection}{\sffamily\normalsize\bfseries}{\thesubsubsection}{0.5em}{}

\pagestyle{fancy}
\fancyhf{}
\fancyhead[L]{\small\sffamily Theseus User Guide}
\fancyhead[R]{\small\sffamily\thepage}
\renewcommand{\headrulewidth}{0.3pt}

\setlength{\parindent}{0pt}

% Convenience commands shared by every guide.
\newcommand{\page}[1]{\textbf{#1}}
\newcommand{\file}[1]{\texttt{#1}}
\newcommand{\route}[1]{\texttt{#1}}
\newcommand{\env}[1]{\texttt{#1}}
\newcommand{\code}[1]{\texttt{#1}}

% Callout boxes used by every guide.
\newtcolorbox{tldr}{
  colback=teal!5!white,colframe=teal!55!black,
  fonttitle=\sffamily\bfseries,title=TL;DR,
  breakable,boxrule=0.4pt,arc=2pt
}
\newtcolorbox{youwillneed}{
  colback=gray!4!white,colframe=gray!55!black,
  fonttitle=\sffamily\bfseries,title=You will need,
  breakable,boxrule=0.4pt,arc=2pt
}
\newtcolorbox{notcovered}{
  colback=orange!5!white,colframe=orange!55!black,
  fonttitle=\sffamily\bfseries,title=This guide does NOT cover,
  breakable,boxrule=0.4pt,arc=2pt
}
\newtcolorbox{preview}{
  colback=yellow!8!white,colframe=yellow!55!black,
  fonttitle=\sffamily\bfseries,title=Preview --- partially shipped,
  breakable,boxrule=0.4pt,arc=2pt
}
\newtcolorbox{nextup}{
  colback=blue!4!white,colframe=blue!55!black,
  fonttitle=\sffamily\bfseries,title=Where to read next,
  breakable,boxrule=0.4pt,arc=2pt
}

% Insert a screenshot only when the file exists, so the build does not
% break before the Playwright capture run is wired up.
\newcommand{\screenshot}[3]{%
  \IfFileExists{screenshots/#1}{%
    \begin{figure}[h]\centering
      \includegraphics[width=\linewidth]{screenshots/#1}
      \caption{#2}\label{fig:#3}
    \end{figure}%
  }{%
    \begin{center}\small\itshape%
      [screenshot \texttt{screenshots/#1} not yet captured --- #2]%
    \end{center}%
  }%
}


\title{\sffamily\bfseries Operator Console\\[2pt]\large Guide 6 of 6}
\date{}
\author{}

\begin{document}
\maketitle
\thispagestyle{empty}

\begin{tldr}
The Operator Console is the founder-only surface of last resort: the
kill switches, the live-authorization toggles, deletion requests,
replication outreach drafts, and the diagnostic loops. None of it is
needed in a normal day. All of it should be reachable in under thirty
seconds when something is wrong. This guide is the map.
\end{tldr}

\begin{youwillneed}
\begin{itemize}[leftmargin=*,itemsep=0.2em,topsep=0.2em]
  \item A signed-in Codex session with operator privileges.
  \item Shell access to the host that runs the launchd
        \code{currents-api} and \code{currents-scheduler} services
        (for the env-var kill switches).
  \item The runbooks under \file{docs/operations/} open in a separate
        tab.
\end{itemize}
\end{youwillneed}

\begin{notcovered}
\begin{itemize}[leftmargin=*,itemsep=0.2em,topsep=0.2em]
  \item How Currents, the Oracle, and Forecasts produce their normal
        output --- Guides 3--5.
  \item Routine triage of principles --- Guide 2.
  \item Security architecture in the abstract; see
        \file{docs/security/Threat\_Model.md}.
\end{itemize}
\end{notcovered}

\section{Map of the console}

The console is not one page. It is a small set of operator-only
routes inside the Codex plus a small set of environment variables on
the host. The routes:

\begin{itemize}[leftmargin=*,itemsep=0.2em]
  \item \route{/(authed)/ops} --- top-level operator surface.
        Scheduler health, ingestion toggles, recent error counts.
  \item \route{/(authed)/admin} --- administrative actions: user /
        org management, deletion requests, retraction workflow.
  \item \route{/(authed)/founder-currents} --- per-opinion controls
        for Currents (revoke, hold, edit event metadata).
  \item \route{/(authed)/forecasts} (with operator role) --- per-row
        \texttt{liveAuthorizedAt} toggles and per-bet confirmation
        modals.
  \item \route{/(authed)/publication} --- the publication review
        queue; \texttt{PublicationReview} actions including unpublish.
  \item \route{/(authed)/social} --- replication outreach drafts and
        manual external posting controls.
\end{itemize}

\screenshot{06_operator_dashboard.png}{The \route{/ops} dashboard, showing scheduler health and the kill-switch state for each ingestion path.}{op-dashboard}

\section{Kill switches}

There is no global kill switch. There are several precise ones, and
that is deliberate. Each disables one specific risky surface.

\begin{itemize}[leftmargin=*,itemsep=0.3em]
  \item \env{CURRENTS\_X\_INGESTION\_DISABLED=true} --- Currents X
        ingestion. Causes the next scheduler cycle to log
        \code{currents.x\_ingestion.disabled reason=manual\_kill\_switch}
        and write no new events. Clearing \env{X\_BEARER\_TOKEN} has
        the same effect.
  \item \env{ARTICLES\_ENABLED=false} --- the long-form article
        dispatcher. Existing articles stay up; no new ones are
        clustered or written.
  \item \env{FORECASTS\_LIVE\_TRADING\_ENABLED=false} --- the master
        live-trading flag from Guide 5. Setting this to false
        disables gate 5 of the eight-gate contract and refuses every
        downstream submit.
  \item \env{FORECASTS\_KILL\_SWITCH=true} --- the per-submit kill
        switch (gate 8). Re-checked immediately before each order
        placement; trips even bets that already passed operator
        confirmation.
\end{itemize}

After flipping any of these, restart the affected service. For the
local launchd Currents runtime, the procedure is documented in
\file{docs/operator/CURRENTS.md} \S``Bringing Currents to life'' and
the helper script
\code{./scripts/refresh-local-currents-runtime.sh --restart --health-check}.

\section{Live authorization, revoked}

The inverse of gates 6--7 in Guide 5. Operator routes:

\begin{itemize}[leftmargin=*,itemsep=0.2em]
  \item \textbf{Revoke per-prediction authorization.} On the
        forecast card, click \page{Revoke live authorization}.
        \texttt{liveAuthorizedAt} is cleared. The row drops out of
        the public grid (the row is publicly visible iff
        \texttt{liveAuthorizedAt IS NOT NULL}).
  \item \textbf{Cancel an open order before fill.} For exchanges that
        permit cancel, the operator UI exposes a cancel button on
        the position row. The submitter logs the cancel; the position
        rolls back.
  \item \textbf{Close an existing position.} The portfolio view has
        a manual close action. This is a real-funds action; it is
        gated behind the same confirmation modal as opening a bet.
\end{itemize}

\section{Deletion requests and retractions}

Two adjacent flows; understand the difference before using them.

\begin{itemize}[leftmargin=*,itemsep=0.3em]
  \item \textbf{Deletion} (\route{/admin}) --- removal of source
        material. Soft-deletes the upload (\texttt{deletedAt} set);
        all derived claims, conclusions, and citations are
        cascade-marked. The public surface stops resolving the
        affected rows. Use for: user / subject deletion requests;
        legal or compliance reasons.
  \item \textbf{Retraction} (\route{/publication}) --- removal of a
        published output. The retraction is itself a published
        artifact: it creates a new
        \texttt{PublishedConclusion} row that retracts the previous
        version, leaving an audit trail. Use for: an article or
        conclusion the firm no longer endorses.
\end{itemize}

A retraction is a stance. A deletion is a removal of evidence. Do not
confuse them.

\section{Replication outreach drafts}

\route{/(authed)/social} contains the operator-drafted outreach
threads that go out to external replicators, journalists, and
collaborators. The page generates a draft from the firm's recent
public material; the operator edits and either posts manually or
exports for sending through a separate channel.

\begin{preview}
Automated sending is not enabled. Today, the outreach surface is a
draft generator plus a send-when-I-say-so review queue. Operator
authorization for automated outreach posting is intentionally not
shipped; see the social outbound roadmap.
\end{preview}

\section{Diagnostic loops}

The fastest way to localize a problem before flipping a switch:

\begin{itemize}[leftmargin=*,itemsep=0.2em]
  \item \code{python noosphere/scripts/diagnose\_currents\_pipeline.py}
        --- the Currents non-writing diagnostic. Booleans and counts
        only; never prints secrets.
  \item \page{OperatorPulse} card on the founder dashboard ---
        Currents ingestion health, color-coded.
  \item \route{/ops} --- recent error counts per scheduler loop.
  \item \file{docs/operator/SCHEDULER\_OPS.md} --- the scheduler
        operator reference.
  \item \file{docs/security/Threat\_Model.md} --- the explicit list
        of which surfaces are allowed to do what.
\end{itemize}

\section{The one-page checklist when something is wrong}

\begin{enumerate}[leftmargin=*,itemsep=0.2em]
  \item Look at the dashboard. Is \page{OperatorPulse} red? What does
        the red text say?
  \item Run the relevant non-writing diagnostic
        (\code{diagnose\_currents\_pipeline.py} for Currents; the
        equivalent forecasts diagnostic for the trading surface).
  \item If a kill switch is the right answer, flip the most specific
        one. Do not disable broader surfaces than you need to.
  \item Restart the affected service with the documented helper
        script; do not stop the entire stack.
  \item Once stable, file a critique through \route{/critiques} so
        the failure has a written record.
\end{enumerate}

\section{Where the docs are}

\begin{itemize}[leftmargin=*,itemsep=0.2em]
  \item \file{docs/operator/CURRENTS.md} --- Currents operator notes.
  \item \file{docs/operator/SCHEDULER\_OPS.md} --- scheduler ops.
  \item \file{docs/operator/public\_surfacing.md} --- public-surface
        predicates.
  \item \file{docs/operations/Forecasts\_Founder\_Alpha\_Runbook.md}
        --- the live-trading runbook.
  \item \file{docs/security/Threat\_Model.md} --- threat model.
\end{itemize}

\begin{nextup}
This is the last of the six guides. After this, the right reading is
the runbooks under \file{docs/operations/}: the scheduler, the
forecasts portfolio setup, and the founder-alpha runbook. The next
scheduled task you should set up is a weekly fifteen-minute slot for
principle triage (Guide 2) and a daily one-minute glance at
\page{OperatorPulse}.
\end{nextup}

\end{document}
